\chapter{Introduction}
\label{introduction}


% ── Background and Motivation ─────────────────────────────────────────────────
\section{Background and Motivation}
\label{sec:intro_motivation}

Training a single large transformer architecture can emit as much carbon as five cars over their entire lifespans \parencite{strubellEnergyPolicyConsiderations2019a}. That figure helped catalyze the Green \gls{ai} movement, a growing body of research that argues computational efficiency must be treated as a primary evaluation criterion alongside predictive accuracy \parencite{schwartzGreenAI2020}. The studies that established this case share a common focus: large-scale deep learning applied to text and images. Yet the majority of deployed machine learning does not look like this. Insurance companies predict churn, banks assess credit risk, particle physicists classify collision events: these are structured, row-column datasets processed by logistic regression, gradient-boosted trees and feedforward networks. Tabular classification is the dominant workload in deployed machine learning. Its ecological footprint has not been systematically measured.

This omission has a concrete practical consequence. A practitioner who wants to make ecologically informed algorithm choices for a tabular task lacks a comprehensive empirical basis to draw on. Whether XGBoost is more efficient than Random Forest on a 30{,}000-row dataset, at what dataset size \gls{gpu} acceleration becomes the greener hardware choice for gradient boosting, or how the cumulative inference cost of a deployed model compares to what was spent training it: these are quantifiable operational questions that remain unanswered.


% ── Problem Statement ─────────────────────────────────────────────────────────
\section{Problem Statement}
\label{sec:intro_problem}

The gap is structural. The Green \gls{ai} literature established that training energy is a measurable and consequential quantity, but empirically it has remained tethered to large-scale deep learning: landmark studies measure transformers trained on text and express costs in \gls{gpu}-hours \parencite{strubellEnergyPolicyConsiderations2019a, dodgeMeasuringCarbonIntensity2022}. The benchmarking literature that rigorously evaluates classical algorithms on tabular data, meanwhile, treats energy as entirely outside its scope and compares models exclusively on predictive accuracy \parencite{grinsztajnWhyTreebasedModels2022a, gorishniyRevisitingDeepLearning2021a}. A systematic study of how energy consumption and predictive performance jointly scale across algorithm families and dataset sizes on tabular data does not exist.

Two further methodological problems widen this gap. First, the software tools commonly used to estimate training emissions have been shown to systematically undercount actual energy consumption \parencite{rodriguezEvaluatingEnergyConsumption2024, fischerGroundTruthingAIEnergy2025}, meaning that published figures for any study that uses them should be treated as lower bounds rather than ground truth. Second, even an accurate measurement of a training run captures only a fraction of a model's true ecological footprint. Hyperparameter search, which for complex models on large datasets can cost more than all benchmark training runs combined, is invisible in published tables \parencite{strubellEnergyPolicyConsiderations2019a}. The inference phase is invisible too, even though repeated prediction over a deployment lifetime has been shown to account for the majority of a production system's cumulative energy demand \parencite{desislavovTrendsAIInference2023}. Consequently, the existing literature often provides an incomplete assessment of the true ecological footprint: the scope is largely restricted to an unrepresentative class of workloads, the tools used systematically undercount the energy they measure, and the training phase alone is treated as a proxy for the full lifecycle cost.


% ── Research Objectives ───────────────────────────────────────────────────────
\section{Research Objectives}
\label{sec:intro_objectives}

This thesis addresses these gaps through an empirical benchmark that treats ecological efficiency as a primary evaluation dimension and extends the measurement scope to cover the full lifecycle of a model. The central research question is:

\begin{researchquestion}
How does the trade-off between energy consumption and predictive performance vary across Logistic Regression, Random Forest, XGBoost and \gls{mlp} models when the complexity of the models and the size of the dataset are systematically varied?
\end{researchquestion}

Five sub-questions structure the investigation:

\begin{itemize}
    \item Under what conditions does the resource consumption of a more complex model exceed the accuracy gain it delivers over a simpler baseline?
    \item How do energy consumption and CO\textsubscript{2} output scale with increasing \gls{mlp} depth and width?
    \item How does energy consumption differ between \gls{cpu}-based and \gls{gpu}-based training for equivalent tasks, and at what dataset size does \gls{gpu} acceleration become ecologically justified?
    \item How does reducing the number of training instances on a large dataset affect model accuracy and energy consumption across different algorithm families?
    \item At what prediction volume does the cumulative inference CO\textsubscript{2} footprint of a deployed model surpass its one-time training footprint, and how does this break-even point differ across models?
\end{itemize}

To answer these questions, the study benchmarks five classifiers on three tabular datasets spanning four orders of magnitude in size: the Wine quality, Default of Credit Card Clients and HIGGS boson detection datasets. The classifiers (Logistic Regression, Random Forest, XGBoost and \gls{mlp}) were selected to represent the full spectrum of tabular classification, from a linear baseline through tree ensembles to a feedforward neural network; XGBoost is evaluated on both \gls{cpu} and \gls{gpu} hardware. A corrected energy measurement methodology addresses the systematic undercounting inherent in standard software-based estimation tools, and these adjusted figures are used throughout. To support cross-model efficiency comparisons that neither raw emissions nor accuracy alone can capture, the study introduces a custom composite metric, the \gls{cf1}, which expresses \gls{co2eq} cost per unit of predictive performance and is defined formally in \Cref{design_metrics}. A counterfactual analysis using one year of hourly grid data further quantifies how much the German electricity grid's diurnal and seasonal variability affects the absolute emissions of the training runs conducted in this study.


% ── Structure of the Thesis ───────────────────────────────────────────────────
\section{Structure of the Thesis}
\label{sec:intro_structure}

The remainder of this thesis is organized as follows. \Cref{theoretical_background} establishes the theoretical foundations: the Green \gls{ai} paradigm, energy and carbon measurement in \gls{ml}, the classification algorithms under study and the statistical evaluation methods. \Cref{related_work} surveys the empirical literature on \gls{ml} energy benchmarking and identifies the structural gap between Green \gls{ai} research and classical tabular classification benchmarks that this work addresses. \Cref{experimental_design} defines the datasets, models, evaluation metrics including \gls{cf1}, validation strategy and hyperparameter tuning protocol. \Cref{implementation} describes the hardware and software environment, the emissions measurement infrastructure and each of the five experiments conducted. \Cref{results} presents the empirical findings across all five experimental dimensions: measurement validation, tuning overhead, the cross-dataset benchmark, hardware and scaling effects, and the lifecycle break-even analysis. \Cref{conclusion} synthesizes the findings, acknowledges the study's limitations and provides actionable implications for Green \gls{ai} practitioners alongside directions for future research.
